\section{Reading Guide and Evidence Map}
\label{sec:archive-guide}

This document is a local research archive accompanying \emph{Evidential Actor Surfaces for Physically Consistent
Driving Reconstruction}. It consolidates the experiments and theory developed from WorldSim V6 through V7.1.
It is intended for private inspection and research continuity, not for OpenReview submission.
The scientific snapshot is the V7.1 branch at commit \texttt{1913ab0e}, including the completed M43 AV2 evaluation.

The main paper follows the current surface--evidence--return formulation. This archive also retains earlier
successful interfaces, rejected alternatives, diagnostic experiments, and historical visualizations.
A positive result retains its original data role: development, reused cohort, independent confirmation,
and synthetic capability are different kinds of evidence.

\begin{table}[h]
\centering\small
\setlength{\tabcolsep}{4pt}
\begin{tabular}{p{.22\linewidth}p{.67\linewidth}}
\toprule
Evidence & Current role\\
\midrule
M7 / M8 & Target-set expansion and frame-supervised canonical geometry; exposed source development.\\
M33--M39 & Producer evidence, separate anchor/child heads, and frozen categorical composition.\\
M22 / M28 & Rigid coordinate composition and physical--visual non-interference.\\
M44 / M49 & Exact CDF responsibility and finite attenuation identities.\\
M50 / M51 & Frame-weighted ray-loss rejection and smooth/minimum operator diagnosis.\\
M43 & Complete frozen AV2 comparison; hit improvement with early-return regression.\\
V6--V7 & Earlier compilation, reliability, and selective-authority results; historical context.\\
\bottomrule
\end{tabular}
\caption{How the retained evidence enters the revised paper.}
\end{table}

\paragraph{Notation and implementation IDs.}
EAS set expansion corresponds to M7, the frame-supervised geometry to M8, and the final categorical
measure to M39 using frozen M35 anchor and M38 child heads.
The unit categorical comparator uses the same M8 centers and scales with unit weights.
The original completion comparator in the geometry table is a point surface scored by the literal operator.
These baselines are distinct. M18 is a learned query decoder; M21 is the earlier frozen unit-energy
candidate; neither denotes the final evidential model.

\paragraph{Corrections to the previous manuscript.}
The revised method distinguishes the alpha-composited expected-depth surrogate from literal
point-surface evaluation. M8's pre-hit term is a hinge on expected depth; M38 separately introduces
integrated pre-hit child optical thickness. M39 freezes heads trained through transmittance and then
changes composition, whereas direct categorical fine-tuning was tested later as M40.
The AV2 cohort is sampled at positions $0,4,\ldots,76$ from the sorted 90-log complement after excluding
60 previously consumed logs.

A read-only implementation audit also identifies a naming collision in M43:
\texttt{\_evaluate\_bundle} first constructs literal point-surface rows, then replaces
\texttt{baseline\_early\_count} and \texttt{baseline\_hit\_count} with the unit categorical counts.
The \texttt{m39\_surface\_return} comparison remains internally matched and supports the reported M43
verdict. However, the \texttt{m8\_point\_surface} early/hit deltas compare a categorical baseline to a
literal output. They cannot be interpreted as a matched geometry-transfer experiment.
Its Chamfer pair is unaffected. We retain the immutable run and omit that mixed-operator causal attribution.

\paragraph{Archive structure.}
The next sections give reproducibility details, complete proofs, and the V6--V6.7 research lineage.
They are followed by the extended V7.1/V7 method and experimental record, plus the restored V7 certificate,
cohort, interval, and camera-evidence appendix.
The separate source snapshot \texttt{supplement\_v7\_legacy.tex} is preserved in the repository.

\clearpage
\setcounter{tocdepth}{2}
\tableofcontents
\clearpage
